\section{Reiteración visible y diferenciación genealógica en una región}
\label{md:ejemplo:region}
Las biografías de \eqref{act21:eq:biografias-tpk-tres-caracteres}
permiten precisar la diferencia entre una reiteración de cifras y un
retorno del estado. Para el canal que se reconoce como \(e\), el
calendario \eqref{act21:eq:calendario-0011-interno} produce
\[
 201101\mid121221\mid102011\mid012222\mid102011.
\]
La elevación precede a la lectura decimal. Las mismas matrices actúan
en los otros dos canales, conservando sus emisiones iniciales y sus
estados regionales distintos. Esta estructura común permite comparar
las biografías sin identificar sus coordenadas.

Sea \(w\) una concatenación de \(L\) trits, y \(N\) su valor entero
en base tres. El cilindro correspondiente es
\[
 I_w=[N/3^L,(N+1)/3^L).
\]
Para \(d\) posiciones decimales, el prefijo es estable en todo el
cilindro exactamente cuando
\[
 \left\lfloor\frac{10^dN}{3^L}\right\rfloor
 =\left\lfloor\frac{10^d(N+1)-1}{3^L}\right\rfloor.
\]
Es la aplicación del lema~\ref{act21:censo:estabilidad} al lector
decimal. El extremo derecho semiabierto explica el término \(-1\)
de la división entera. Añadiendo la parte entera de la coordenatización regional,
las cinco lecturas sucesivas son
\begin{center}
\begin{tabular}{@{}rl@{}}
\toprule
Longitud trítica & Prefijo estable de la lectura regional\\
\midrule
6 & \(2.71\ldots\)\\
12 & \(2.71828\ldots\)\\
18 & \(2.7182818\ldots\)\\
24 & \(2.7182818284\ldots\)\\
30 & \(2.71828182845904\ldots\)\\
\bottomrule
\end{tabular}
\end{center}
En longitud veinticuatro se tiene
\(N=202864003874\) y \(3^{24}=282429536481\).
En longitud treinta se obtiene
\(N=147887858824447\) y \(3^{30}=205891132094649\).
Las divisiones enteras anteriores certifican los prefijos de la tabla.
En particular, las primeras doce cifras fraccionarias se agrupan como
\[
 718281828459
 =7\,\underbrace{1828\,1828}_{\text{reiteración posicional}}\,459.
\]
El cilindro de veinticuatro trits ya fija ambas copias y la cifra
siguiente. Esta localización distingue el episodio decimal de la
reaparición del bloque trítico \(102011\), situado en las posiciones
tercera y quinta: sus cortes y sus bases de lectura son diferentes.

La repetición del bloque visible tampoco identifica sus preestados.
Los preestados módulo veintisiete conservados en el desarrollo de
esta biografía son
\[
 (25,11,7,17,8,16),\qquad (7,8,4,23,26,7),
\]
con cocientes posteriores respectivos
\[
 (6,13,9,2,13,1),\qquad (15,0,3,19,23,25).
\]
La igualdad de la emisión coexiste así con una genealogía diferenciada.
Su función explicativa es concreta: la observación de un bloque
reiterado conserva menos información que la continuación enriquecida
que lo produce. Para atribuir una diferencia futura exclusivamente
a la memoria hay que conservar asimismo las fases y las coordenatizaciones de
comparación.

Este episodio relaciona tres niveles: las transiciones producen la
cadena trítica, los cilindros determinan las cifras estables y el
estado conserva la procedencia de una emisión reiterada. La
profundidad arbitraria pertenece a la prolongación demostrada
anteriormente; esta lectura finita precisa un episodio de esa
biografía. La regularidad que interesa estudiar reside en la
compatibilidad entre los niveles, no en prolongar indefinidamente
un fragmento decimal: una tercera copia de \(1828\) comenzaría por
\(1\), mientras la continuación ya producida comienza por \(4\).
